\clearpage
\VerdierPhysicalPageAnchor{0163}%
\begingroup\small\setlength{\parskip}{0pt}\setlength{\emergencystretch}{1.5em}
\VerdierSemanticAnchor{ch4-4-1-2-b}
b) $\mathcal{F}\ell_2(\mathcal J)$ から $\mathcal D$ の完全三角の圏への、
次の形の関手
\[
\omega_2=(\mathord\bullet\xrightarrow f\mathord\bullet\xrightarrow g\mathord\bullet)
\longmapsto
X_f\xrightarrow{X*u_1(\omega_2)}X_{gf}
\xrightarrow{X*u_2(\omega_2)}X_g
\xrightarrow{\delta(\omega_2)}X_f[1].
\]
型 $\mathcal J$ のスペクトル対象を表すため、記号
\[
(X_f, f\in\operatorname{Ob}\mathcal{F}\ell_1(\mathcal J),\
\delta(\omega_2),
\ \omega_2\in\operatorname{Ob}\mathcal{F}\ell_2(\mathcal J)),
\]
を用いる。混同のおそれがなければ、単に $(X_f,\delta)$ と書く。

\VerdierSemanticAnchor{ch4-4-1-3}
\medskip\noindent\textbf{4.1.3.}\quad
$C$ を加法圏とし、
$F\colon\mathcal J\longrightarrow\operatorname{comp}(C)$ を関手とする。
第 I 章 (3.3.6) で、$F$ に型 $\mathcal J$ で $K(C)$ に値をとる
スペクトル対象を付随させる方法を示した。第 I 章 (3.3.6) で述べた
スペクトル対象を、\emph{関手 $F$ に付随するスペクトル対象}という。

\VerdierSemanticAnchor{ch4-4-1-4}
\medskip\noindent\textbf{4.1.4.}\quad
$\mathcal A$ をアーベル圏、$\mathcal J$ を圏とする。
\emph{型 $\mathcal J$ で $\mathcal A$ に値をとるスペクトル対象}は、
次のものからなる：

a) 関手列
$H^n\colon\mathcal{F}\ell_1(\mathcal J)\longrightarrow\mathcal A$、
$n\in\mathbb Z$。

b) 各対象
$\omega_2=(\mathord\bullet\xrightarrow f\mathord\bullet\xrightarrow g\mathord\bullet)$
を $\mathcal{F}\ell_2(\mathcal J)$ の対象とし、各整数 $n$ に対し、射
$\delta^n(\omega_2)\colon H^n(g)\to H^{n+1}(f)$ が与えられる。この射は
$\omega_2$ に関して関手的であり、図式
\[
\cdots\longrightarrow H^n(f)
\xrightarrow{H^n*u_1(\omega_2)}H^n(gf)
\xrightarrow{H^n*u_2(\omega_2)}H^n(g)
\xrightarrow{\delta^n(\omega_2)}H^{n+1}(f)
\longrightarrow\cdots
\]
が $\mathcal A$ の両側無限完全列となる。

型 $\mathcal J$ のスペクトル対象を表すため、記号
\[
(H^n(f), f\in\operatorname{Ob}\mathcal{F}\ell_1(\mathcal J),\
\delta^n(\omega_2),
\ \omega_2\in\operatorname{Ob}\mathcal{F}\ell_2(\mathcal J))
\]
を用いる。混同のおそれがなければ、単に $(H^n(f),\delta)$ と書く。

\VerdierSemanticAnchor{ch4-4-1-5}
\medskip\noindent\textbf{4.1.5.}\quad
$F\colon\mathcal D\to\mathcal D'$ を二つの三角圏の間の完全関手とし、
$(X_f,\delta)$ を型 $\mathcal J$ で $\mathcal D$ に値をとる
スペクトル対象とする。関手の対 $(FX_f,F\delta)$ は、型 $\mathcal J$ で
$\mathcal D'$ に値をとるスペクトル対象である。同様に、三角圏から
アーベル圏へのコホモロジー的関手
$H\colon\mathcal D\to\mathcal A$ は、$\mathcal D$ に値をとる
スペクトル対象を $\mathcal A$ に値をとるスペクトル対象に移す。

\VerdierSemanticAnchor{ch4-4-1-6}
\medskip\noindent\textbf{4.1.6.}\quad
$\mathcal D$ を三角圏とし、$(X_f,\delta)$、$(Y_f,\delta')$ を型
$\mathcal J$ で $\mathcal D$ に値をとる二つのスペクトル対象とする。
\emph{スペクトル対象の射}とは、関手の射
\[
m_f\colon X_f\longrightarrow Y_f,\qquad
f\in\operatorname{Ob}\mathcal{F}\ell_1(\mathcal J),
\]
\endgroup
